% Official solution ("Problem 3 B. Marking scheme -- Observer on an extrasolar
% star, 25 points", 2014_da_sol.pdf pp.36-39). Decimal commas kept as printed.
% Note: the official solution names the asterism "Big Dipper" (and keeps some
% Romanian words: Planeta, Soare, Pamant); the official English statement --
% and the book's problem text -- uses the Mizar system instead. Transcribed
% as printed.

\begin{description}
\item[a)] The magnitude of the \textbf{SUN as seen from the observer on the
  Sirrius planet} % sic: "Sirrius" in source
  \[ \log\frac{E_{\mathrm{SUN},EARTH}}{E_{\mathrm{SUN},PLANET}}
     = -0{,}4\left(m_{\mathrm{SUN},EARTH} - m_{\mathrm{SUN},PLANET}\right), \]
  Where: $E$ represents the apparent brightness of the body at the
  corresponding distance to the observer.
  \begin{align*}
  \log\frac{\dfrac{L_{\mathrm{S}}}{4\pi d_{\mathrm{S},E}^2}}
           {\dfrac{L_{\mathrm{S}}}{4\pi d_{\mathrm{S,P}}^2}}
    &= -0{,}4\left(m_{\mathrm{S},E} - m_{\mathrm{S,P}}\right);\\
  d_{\mathrm{S},E} = 1\,\mathrm{A}U;\
  d_{\mathrm{S,P}} &= 2{,}6\,\mathrm{pc}
    = 2{,}6\cdot206265\ \mathrm{A}U;\\
  m_{\mathrm{S},E} &= -26{,}78^{\mathrm{m}};\\
  m_{\mathrm{S,P}} &= 1{,}87^{\mathrm{m}},
  \end{align*}
  Apparent magnitude of Sun seen from the observer on one of the Sirius
  planet.
\item[b)] \textbf{Magnitude of Sirius, seen from the observer from its
  planet}:
  \begin{align*}
  \log\frac{E_{\mathrm{Sirius,Planeta}}}{E_{\mathrm{Sun,Planeta}}}
    &= -0{,}4\left(m_{\mathrm{Sirius,Planeta}}
       - m_{\mathrm{Sun,Planeta}}\right),\\
  d_{\mathrm{Sun,Planeta}} &= 2{,}6\,\mathrm{pc}
    = 2{,}6\cdot206265\ A\mathrm{U};\\
  d_{\mathrm{Sirius,Planeta}} &= 10\ A\mathrm{U};\\
  L_{\mathrm{Sirius}} &= 25\cdot L_{SUN};\\
  m_{\mathrm{Sun,Planeta}} &= 1{,}87^{\mathrm{m}};\\
  m_{\mathrm{Sirius,Planeta}} &= -25{,}255^{\mathrm{m}},
  \end{align*}
\item[c)] The luminosity of \textbf{Big Dipper}:
  \[ L_{\mathrm{Big\ Dipper}} = L_1 + L_2 + L_3 + L_4 + L_5 + L_6 + L_7, \]
  where $L_1$, $L_2$, \dots, $L_7$ are the luminosties % sic: "luminosties"
  of the seven stars of the asterism:
  \begin{align*}
  \log\frac{E_{\mathrm{S}un,Earth}}{E_{1,Earth}}
    &= -0{,}4\left(m_{\mathrm{SUN},Earth} - m_{1,Earth}\right),\\
  \log\frac{\dfrac{L_{\mathrm{S}un}}{4\pi d_{\mathrm{S}un,Earth}^2}}
           {\dfrac{L_1}{4\pi d_{1,Earth}^2}}
    &= -0{,}4\left(m_{\mathrm{S}un,Earth} - m_{1,Earth}\right);\\
  d_{1,Earth} &= 124\cdot63240\ A\mathrm{U};\\
  m_{\mathrm{S}un,Earth} = -26{,}78^{\mathrm{m}};\
  m_{1,Earth} &= 1{,}8^{\mathrm{m}};\\
  L_1 &\approx 225\cdot L_{\mathrm{S}},
  \end{align*}
  The luminosity of star 1 (Dubhe) from the asterism Big Dipper.

  Symilarly for the rest of stars in the asterism: % sic: "Symilarlyfor"
  \begin{align*}
  L_2 &\approx 50\cdot L_{\mathrm{S}};\\
  L_3 &\approx 60\cdot L_{\mathrm{S}};\\
  L_4 &\approx 0{,}134\cdot L_{\mathrm{S}};
    % sic: 0,134 as printed; the data of the problem (m = 3,3 at 58 ly)
    % give L_4 of order 13 L_S
  \\
  L_5 &\approx 52\cdot L_{\mathrm{S}};\\
  L_6 &\approx 4\cdot L_{\mathrm{S}};\\
  L_7 &\approx 1{,}37\cdot L_{\mathrm{S}}.
  \end{align*}
  The luminosity of the asterism Big Dipper is:
  \begin{align*}
  L_{\mathrm{Big\ Dipper}} &= L_1 + L_2 + L_3 + L_4 + L_5 + L_6 + L_7;\\
  L_{\mathrm{Big\ Dipper}} &= (225 + 50 + 60 + 0{,}134 + 52 + 4
    + 1{,}37)L_{Sun};\\
  L_{\mathrm{Big\ Dipper}} &= 392{,}5\cdot L_{\mathrm{S}un}.
  \end{align*}
\item[d)] \textbf{The distance between Big Dipper and Earth}
  $d_{\mathrm{Big\ Dipper},Earth}$:
  \begin{align*}
  \log\frac{E_{\mathrm{Big\ Dipper},Earth}}{E_{\mathrm{Soare},Earth}}
    &= -0{,}4\left(m_{\mathrm{Big\ Dipper},Earth}
       - m_{\mathrm{Soare},Earth}\right),\\
  L_{\mathrm{Big\ Dipper}} &= 392{,}5\cdot L_{\mathrm{S}};\\
  d_{\mathrm{Big\ Dipper},Earth} &= 573\cdot10^{3}\cdot\sqrt{392{,}5}
    \cdot d_{\mathrm{Soare},Earth};\\
  d_{\mathrm{Big\ Dipper},Earth} &\approx 11{,}3\cdot10^{6}\ AU.
  \end{align*}
\item[e)] Geocentric angular distance $\Delta\theta$, \textbf{between
  Big Dipper and Sirius}. % sic: source prints "between dintre Big Dipper"

  The equatorial coordinates of asterism Big Dipper $(\sigma_1)$ and Sirius
  $(\sigma_2)$ located on the sphere of geocentric celestial sphere are:
  \begin{align*}
  \alpha_{\mathrm{Big\ Dipper}} &= \alpha_1
    = 13^{\mathrm{h}}23^{\mathrm{min}}55{,}5^{\mathrm{s}};\
  \delta_{\mathrm{Big\ Dipper}} = \delta_1 = 54^\circ55'31'';\\
  \alpha_{\mathrm{Sirius}} &= \alpha_2 = 6^{\mathrm{h}}45^{\mathrm{min}};\
  \delta_{\mathrm{Sirius}} = \delta_2 = -16^\circ43'.
  \end{align*}
  % FIGURE NEEDED: celestial-sphere diagram (2014_da_sol.pdf p.38): sphere
  % with poles P, P', equator Q'Q, vernal point gamma, the two stars
  % sigma_1(alpha_1; delta_1) (Big Dipper, point B) and
  % sigma_2(alpha_2; delta_2) (Sirius, point C), their hour-circle feet
  % sigma_01, sigma_02, the angle Delta-theta at O and the right angle at A.

  In the spherical dreptunghic triangle ABC: % sic: "dreptunghic"
  % (Romanian: right-angled)
  \[ \sphericalangle\left(\sigma_1\,\mathrm{O}\,\sigma_2\right)
     = \Delta\theta. \]
  Also
  \begin{align*}
  \sphericalangle\left(\gamma\,\mathrm{O}\,\sigma_{01}\right)
    &= \alpha_{\mathrm{Big\ Dipper}} = \alpha_1;\
  \sphericalangle\left(\sigma_{01}\,\mathrm{OB}\right)
    = \delta_{\mathrm{Big\ Dipper}} = \delta_1;\\
  \sphericalangle\left(\gamma\,\mathrm{O}\,\sigma_{02}\right)
    &= \alpha_{\mathrm{Sirius}} = \alpha_2;\
  \sphericalangle\left(\sigma_{02}\,\mathrm{OC}\right)
    = \delta_{\mathrm{Sirius}} = \delta_2;\\
  \sphericalangle\left(\sigma_{01}\,\mathrm{O}\,\sigma_{02}\right)
    &= \sphericalangle\left(\mathrm{AOC}\right)
    = \alpha_1 - \alpha_2 = \Delta\alpha;\\
  \sphericalangle\left(\sigma_1\,\mathrm{OA}\right)
    &= \sphericalangle\left(\sigma_{01}\,\mathrm{O}\,\sigma_1\right)
     + \sphericalangle\left(\sigma_{01}\,\mathrm{OA}\right)
    = \delta_{\mathrm{Big\ Dipper}} - \delta_{\mathrm{Sirius}}
    = \Delta\delta;\\
  \sphericalangle\left(\sigma_1\,\mathrm{O}\,\sigma_2\right)
    &= \Delta\theta.
  \end{align*}
  Results
  \begin{align*}
  \cos a &= \cos b\cdot\cos c + \sin b\cdot\sin c\cdot\cos A;\\
  A &= 90^{0};\\
  \cos a &= \cos b\cdot\cos c;\\
  a = \Delta\theta;\ b &= \Delta\alpha;\quad c = \Delta\delta;\\
  \cos\Delta\theta &= \cos\Delta\alpha\cdot\cos\Delta\delta;\\
  \Delta\alpha = \alpha_1 - \alpha_2
    &= 13^{\mathrm{h}}23^{\mathrm{min}}55{,}5^{\mathrm{s}}
     - 6^{\mathrm{h}}45^{\mathrm{min}}
     = 6^{\mathrm{h}}38^{\mathrm{min}}55{,}5' \approx 100^\circ;
    % sic: 55,5' for seconds as printed
  \\
  \Delta\delta = \delta_1 - \delta_2
    &= 54^\circ55'31'' + 16^\circ43' = 71^\circ38'31''
    \approx 72{,}4^\circ;\\
  \cos\Delta\theta &= \cos\left(100^\circ\right)
    \cdot\cos\left(72{,}4^\circ\right)
    = \left(-0{,}173648177\right)\cdot\left(0{,}30236989\right)
    = -0{,}05250598;\\
  \Delta\theta &\approx 90{,}3^\circ \approx 90^\circ,
  \end{align*}
\item[f)] The distance between \textbf{Big Dipper and the observer on sirius
  planet of Sirius}, % sic: as printed
  $d_{\mathrm{Big\ Dipper,Planeta}}$:
  \begin{align*}
  d_{\mathrm{Big\ Dipper,Planeta}}
    &\approx d_{\mathrm{Big\ Dipper,Sirius}};\\
  d_{\mathrm{Sirius},Earth} &\approx d_{\mathrm{Sirius,Soare}}
    \approx d_{\mathrm{Soare,Planeta}};\\
  d_{\mathrm{Sirius,Pamant}} &= 2{,}6\,\mathrm{pc};\\
  d_{\mathrm{Soare,Planeta}} &= 2{,}6\,\mathrm{pc}
    = 2{,}6\cdot206265\ \mathrm{UA} = d_{Earth};\\
  d_{\mathrm{Big\ Dipper},Earth} &\approx 11{,}3\cdot10^{6}\ AU;
  \end{align*}
  % FIGURE NEEDED: triangle diagram (2014_da_sol.pdf p.39): right triangle
  % with "Earth + Sun" at the right-angle vertex, "Sirius + Planeta" above
  % it at distance d_Sirius,Pamant, "Big Dipper" to the right at distance
  % d_BigDipper,Earth, hypotenuse d_BigDipper,Sirius, and the angle
  % Delta-theta ~ 90 degrees marked at Earth.
  \[ d_{\mathrm{Big\ Dipper,Sirius}}
     = \sqrt{d_{\mathrm{Big\ Dipper,Pamant}}^2
       + d_{\mathrm{Sirius},Eaarth}^2} % sic: "Eaarth" in source
     = d_{\mathrm{Big\ Dipper,Planeta}}
     \approx 11{,}312\cdot10^{6}\ \mathrm{UA}. \]
\item[g)] The magnitude of \textbf{Big Dipper for an observer from the
  observer on planet of Sirius}: % sic: as printed
  \begin{align*}
  \log\frac{E_{\mathrm{Big\ Dipper,Planeta}}}
           {E_{\mathrm{Big\ Dipper},Earth}}
    &= -0{,}4\left(m_{\mathrm{Big\ Dipper,Planeta}}
       - m_{\mathrm{Big\ Dipper},Earth}\right),\\
  m_{\mathrm{Big\ Dipper},Earth} \approx 2^{\mathrm{m}};\
  d_{\mathrm{Big\ Dipper},Earth} &\approx 11{,}3\cdot10^{6}\ A\mathrm{U};\\
  d_{\mathrm{Big\ Dipper,Sirius}} = d_{\mathrm{Big\ Dipper,Planeta}}
    &\approx 11{,}312\cdot10^{6}\ \mathrm{A}U;\\
  m_{\mathrm{Big\ Dipper,Planeta}}
    &= 2^{\mathrm{m}} + 5^{\mathrm{m}}\cdot
       \log\frac{11{,}312\cdot10^{6}\ \mathrm{AU}}
                {11{,}3\cdot10^{6}\ \mathrm{AU}};\\
  m_{\mathrm{Big\ Dipper,Planeta}}
    &= 2^{\mathrm{m}} + 5^{\mathrm{m}}\cdot
       \log\frac{11{,}312}{11{,}3};\\
  m_{\mathrm{Big\ Dipper,Planeta}} &\approx 2{,}0023^{\mathrm{m}},
  \end{align*}
\end{description}
